\chapter{2013年上海卷（秋理）}

\section{填空题}

\begin{problem}
计算: \(\lim_{n\to \infty}\frac{n + 20}{3n + 13} =\)
\end{problem}

\begin{problem}
设 \(m \in \R\)，\(m^2 + m - 2 + (m^2 - 1)\) 是纯虚数，其中 i 是虚数单位，则 \(m = \)  \fillinblank{}.
\end{problem}

\begin{problem}
若 \(\left|\frac{x^2}{-1} - 1\right|^2 = \left|\frac{x}{y} - 1\right|\)，则 \(x + y = \)  \fillinblank{}.
\end{problem}

\begin{problem}
已知 \(\triangle ABC\) 的内角 \(A, B, C\) 所对的边分别为 \(a, b, c\). 若 \(3a^2 + 2ab + 3b^2 - 3c^2 = 0\)，则角 \(C\) 的大小是 \fillinblank{}. (结果用二元函数值表示)
\end{problem}

\begin{problem}
设常数 \(a \in \R\). 若 \(\left(x^2 + \frac{a}{x}\right)^5\) 的二项展开式中 \(x^2\) 项的系数为 \(-10\)，则 \(a = \)  \fillinblank{}.
\end{problem}

\begin{problem}
方程 \(\frac{3}{3^x - 1} + \frac{1}{3} = 3^{x - 1}\) 的实数为 \fillinblank{} \fillinblank{}.
\end{problem}

\begin{problem}
在极坐标系中，曲线 \( \rho = \cos\theta + 1 \) 与 \( \rho \cos\theta = 1 \) 的公共点到极点的距离为 \fillinblank{} \fillinblank{}.
\end{problem}

\begin{problem}
盒子中装有编号为 1, 2, 3, 4, 5, 6, 7, 8, 9 的九个球，从中任意取出两个，则这两个球的编号之积为偶数的概率是 \fillinblank{}. （结果用最简分数表示）
\end{problem}

\begin{problem}
设 AB 是椭圆 T 的长轴，点 C 在 T 上，且 \( \angle CBA = \frac{T}{4} \) ，若 AB = 4, BC = \( \sqrt{2} \) ，则 T 的两个焦点之间的距离为 \fillinblank{} \fillinblank{}.
\end{problem}

\begin{problem}
设非零常数 \(d\) 是等差数列 \(x_{1}, x_{2}, x_{3}, \cdots, x_{10}\) 的公差，随机变量 \(\xi\) 等可能地取值 \(x_{1}, x_{2}, x_{3}, \cdots, x_{10}\)，则方差 \(D\xi = \fillinblank{}\).
\end{problem}

\begin{problem}
若 \(\cos x\cos y + \sin x\sin y = \frac{1}{2}\)，\(\sin 2x + \sin 2y = \frac{3}{2}\)，则 \(\sin (x + y) =\)  \fillinblank{}.
\end{problem}

\begin{problem}
设 \(a\) 为实常数，\(y = f(x)\) 是定义在 \(\R\) 上的奇函数，当 \(x < 0\) 时，\(f(x) = 0x + \frac{a^2}{x} + 7\)，若 \(f(x) \geqslant a + 1\) 对一切 \(x \geqslant 0\) 成立，则 \(a\) 的取值范围为 \fillinblank{} \fillinblank{}.
\end{problem}

\begin{problem}
在 \(xOy\) 平面上，将两个半圆弧 \((x - 1)^2 + y^2 = 1 (x \geqslant 1)\) 和 \((x - 3)^2 + y^2 = 1 (x \geqslant 3)\)，两条直线 \(y = 1\) 和 \(y = -1\) 围成的封闭图形记为 \(D\)，如图中阴影部分. 记 \(D\) 绕 \(y\) 轴旋转一周而成的几何体为 \(\Omega\)，过 \((0, y) (|y| \leqslant 1)\) 作 \(\Omega\) 的水平截面，所得截面面积为 \(4\pi \sqrt{1 - y^2} + 8\pi\)，试利用祖暅原理，一个平放的圆柱和一个长方体，得出 \(\Omega\) 的体积值为 \fillinblank{} \fillinblank{}.
\end{problem}

\begin{problem}
对区间 \(I\) 上有定义的函数 \(g(x)\)，记 \(g(I) = \{y \mid y = g(x), x \in I\}\)，已知定义域为 \([0,3]\) 的函数 \(y = f(x)\) 有反函数 \(y = f^{-1}(x)\)，且 \(f^{-1}([0,1]) = [1,2)\)，\(f^{-1}((2,4]) = [0,1)\)，若方程 \(f(x) - x = 0\) 有解 \(x_0\)，则 \(x_0 = \fillinblank{}\).
\end{problem}

\section{单选题}

\begin{problem}
设常数 \(a \in \R\)，集合 \(A = \{x \mid (x - 1)(x - a) \geqslant 0\}\)，\(B = \{x \mid x \geqslant a - 1\}\). 若 \(A \cup B = \R\)，则 \(a\) 的取值范围为（\quad）

\choices
    {\((-\infty, 2)\)}
    {\((-\infty, 2]\)}
    {\((2, +\infty)\)}
    {\([2, +\infty)\)}
\end{problem}

\begin{problem}
钱大姐常说“便宜没好货”，她这句话的意思是：“不便宜”是“好货”的（\quad）

\choices
    {充分条件}
    {必要条件}
    {充分必要条件}
    {既非充分也非必要条件}
\end{problem}

\begin{problem}
在数列 \(\{a_{n}\}\) 中，\(a_{n} = 2^{n} - 1\)，若一个 7 行 12 列的矩阵的第 i 行第 j 列的元素 \(a_{ij} = a_{i} \cdot a_{j} + a_{i} + a_{j}\) (\(i = 1, 2, \cdots, 7; j = 1, 2, \cdots, 12\)) 则该矩阵元素能取到的不同数值的个数为（\quad）

\choices
    {18}
    {28}
    {48}
    {63}
\end{problem}

\begin{problem}
在边长为 1 的正六边形 ABCDEF 中，记以 A 为起点，其余顶点为终点的向量分别为 \( \overrightarrow{a_{1}}, \overrightarrow{a_{2}}, \overrightarrow{a_{3}}, \overrightarrow{a_{4}}, \overrightarrow{a_{5}} \) ; 以 D 为起点，其余顶点为终点的向量分别为 \( \overrightarrow{d_{1}}, \overrightarrow{d_{2}}, \overrightarrow{d_{3}}, \overrightarrow{d_{4}}, \overrightarrow{d_{5}} \) . 若 m, M 分别为 \( (\overrightarrow{a_{1}} + \overrightarrow{a_{2}} + \overrightarrow{a_{3}}) \cdot (\overrightarrow{d_{4}} + \overrightarrow{d_{5}} + \overrightarrow{d_{6}}) \) 的最小值、最大值，其中 \( \{i, j, k\} \subseteq \{1, 2, 3, 4, 5\} \) ，\( \{r, s, t\} \subseteq \{1, 2, 3, 4, 5\} \) ，则 m、M 满足（\quad）

\choices
    {\(m = 0, M > 0\)}
    {\(m < 0,M > 0\)}
    {\(m < 0,M = 0\)}
    {\(m < 0,M < 0\)}
\end{problem}

\section{解答题}

\begin{problem}
如图，在长方体 \(ABCD - A'B'C'D'\) 中，\(AB = 2\)，\(AD = 1\)，\(AA' = 1\)，证明直线 \(BC'\) 平行于平面 \(D'AC\)，并求直线 \(BC'\) 平到平面 \(D'AC\) 的距离.
\end{problem}

\begin{problem}
甲厂以 \(x\) 千克/小时的速度运输生产某种产品（生产条件要求 \(1 \leqslant x \leqslant 10\)），每一小时可获得利润是 \(100\left(5x + 1 - \frac{3}{2}\right)\) 元. (1)要使生产该产品2小时获得的利润不低于3000元，求x的取值范围; (2)要使生产900千克该产品获得的利润最大，问:甲厂应该选取何种生产速度\(\perp\) 并求最大利润.
\end{problem}

\begin{problem}
已知函数 \( f(x) = 2\sin \omega x \). 其中常数 \( \omega > 0 \). (1) 若 \( y = f(x) \) 在 \(\left[-\frac{\pi}{4}, \frac{2\pi}{3}\right]\) 上单调递增，求 \(\omega\) 的取值范围; (2) 令 \(\omega = 2\)，将函数 \(y = f(x)\) 的图象向左平移 \(\frac{\pi}{6}\) 个单位，再向上平移 1 个单位，得到函数 \(y = g(x)\) 的图象. 区间 \([a, b]\) (\(a, b \in \R\) 且 \(a < b\)) 满足：\(y = g(x)\) 在 \([a, b]\) 上至少含有 30 个零点. 在所有满足上述条件的 \([a, b]\) 中，求 \(b - a\) 的最小值. 
\end{problem}

\begin{problem}
如图，已知双曲线 \(C: \frac{x^2}{2} - y^2 = 1\)，曲线 \(C_2: |y| = |x| + 1\). \(P\) 是平面内一点，若存在过点 \(P\) 的直线与 \(C_1, C_2\) 都有公共点，则称 \(P\) 为 \(C_1 - C_2\) 型点 (1) 在正确证明 \( C_{1} \) 的左焦点是“ \( C_{1}-C_{2} \) 型点”时，要使用一条过该焦点的直线，试写出一条这样的直线的方程 (不要求验证); (2) 设直线 \( y = kx \) 与 \( C_2 \) 有公共点. 求证 \( |k| > 1 \)，进而证明原点不是“\( C_1 - C_2 \) 型点”; (3) 求证: 圆 \(x^{2} + y^{2} = \frac{1}{2}\) 内的点都不是“\(C_{1} - C_{2}\) 型点”.
\end{problem}

\begin{problem}
给定常数 \(c > 0\)，定义函数 \(f(x) = 2|x + c + 4| - |x + c|\)，数列 \(a_1, a_2, a_3, \ldots\) 满足 \(a_{n+1} = f(a_n), n \in \mathbf{N}^*\). (1) 若 \(a_{1} = -c - 2\)，求 \(a_{2}\) 及 \(a_{3}\); (2) 求证：对任意 \(n \in \mathbf{N}^n, a_{n+1} - a_n \geqslant c\); (3) 是否存在 \(a_{1}\)，使得 \(a_{1}, a_{2}, \cdots, a_{n} \cdots\) 成等差数列\(\perp\) 若存在，求出所有这样的 \(a_{1}\)，若不存在，说明理由.
\end{problem}

